\documentclass[10pt,a4paper]{amsart}
\usepackage[margin=19mm]{geometry}
\usepackage{amsmath,amssymb,mathtools}
\usepackage[scr=rsfs,cal=euler]{mathalfa}
\usepackage{xfrac}
\usepackage[unicode,hidelinks,bookmarks=false]{hyperref}
\newcommand{\ie}{i.e.\ }
\newcommand{\cf}{cf.\ }
\newcommand{\resp}{respectively\ }
\newcommand{\Subsection}{\subsection}
\providecommand{\nobreakdash}{\nobreak\hbox{-}}
\setlength{\emergencystretch}{2em}
\multlinegap0pt
\usepackage[shortlabels]{enumitem}
\usepackage{amssymb}
\usepackage{mathtools}
\usepackage{bm}
\usepackage{float}
\usepackage{xcolor}
\newtheorem{hypothesis}{Hypothesis}
\newtheorem{definition}{Definition}
\newtheorem{lemma}{Lemma}
\newtheorem{corollary}{Corollary}
\newtheorem{theorem}{Theorem}
\newtheorem{proposition}{Proposition}
\newcommand{\Cz}{\mathcal{C}_0}
\newcommand{\Fclass}[2]{\mathcal{F}^{(#1)}_{#2}}
\newcommand{\Kerr}{\mathrm{Kerr}}
\newcommand{\Sym}{\mathrm{Sym}}
\newcommand{\ordc}{\mathrm{ord}_c}
\title{\textbf{Computational Superiority of Non-Markovian Kerr Feedback\\ in Continuous-Variable Quantum Reservoir Computing}}
\author{Daniel Soh}
\date{Source: arXiv:2606.06689v1 (CC BY 4.0)}
\begin{document}
\maketitle

\noindent\emph{Source and licence.} \textbf{Computational Superiority of Non-Markovian Kerr Feedback\\ in Continuous-Variable Quantum Reservoir Computing}. Version arXiv:2606.06689v1 (CC BY 4.0), distributed by arXiv under the Creative Commons Attribution 4.0 International licence, http://creativecommons.org/licenses/by/4.0/, as recorded in the arXiv licence metadata for this exact version and shown on the abstract page https://arxiv.org/abs/2606.06689v1. Full licence notice: \texttt{licenses/arxiv-2606.06689-CC-BY-4.0.txt}.\par\smallskip

\noindent\textit{Editorial scope.} Counted unit: the article's proposition prop:degree (source0.tex lines 1249--1258). The article's complete original proof of it is delivered with it (source0.tex lines 1260--1297, one \texttt{proof} environment of 38 lines). No mathematics is written, repaired or re-derived: the statement and the proof are the article's own lines, in the article's own notation, and no retained line is substituted or re-flowed.  Retained uncounted as the source-local context the proof needs: lines 293--297; lines 313--318; lines 320--325; lines 331--334; lines 353--358; lines 411--416; lines 424--443; lines 473--491; lines 589--600; lines 1098--1130.  Nothing was omitted from any retained span, so the package carries no omission record.  No retained line contains a citation, so the wrapper carries no bibliography and no bibliography entry.  No retained line contains a figure environment, an image-inclusion command or drawing code, so no image is delivered and the package declares no asset dependency.  Editorial formatting only, in addition: an AMS \texttt{amsart} preamble that restates the article's own declarations and macros used by the retained lines, namely the environments \texttt{hypothesis}, \texttt{definition}, \texttt{lemma}, \texttt{corollary}, \texttt{theorem}, \texttt{proposition} on separate counters, together with the standard packages those lines need.  The retained environments print in this document as Hypothesis 1, Definition 1, Lemma 1, Corollary 1, Theorem 1, Proposition 1 in order of appearance, whereas the article prints the same items with the numbering of its own full document.  The complete native source of the article as distributed in the arXiv e-print for this exact version is bundled unchanged as \texttt{sources/202205/source0.tex}.
\begin{equation}
\beta:\mathbb{R}\to\mathbb{R}^{2N},\qquad
\gamma:\mathbb{R}\to\Sym^2(\mathbb{R}^{2N}),
\label{eq:inject}
\end{equation}

\begin{hypothesis}[Injected encoding]\label{H1}
The input enters only through the injected ancilla data \eqref{eq:inject}.
Consequently the per-step reservoir propagation (mirror coupling, the MZI linear
network, intrinsic loss, and the ancilla reset) is a fixed, input-independent
Gaussian channel $\Cz$.
\end{hypothesis}

\begin{hypothesis}[Echo state / contraction]\label{H2}
The single-step covariance propagator $S_1$ induced by $\Cz$ on
$\mathbb{R}^{2N}$ has spectral radius $\rho(S_1)<1$. Equivalently the
continuous-time drift $A$ is Hurwitz, $\operatorname{Re}\lambda<0$ for every
eigenvalue $\lambda$ of $A$.
\end{hypothesis}

\begin{hypothesis}[Active squeezing]\label{H4}
The covariance injection is nonconstant, $\gamma_1\neq0$, and the encoding is
non-degenerate on the measured channel.
\end{hypothesis}

\begin{equation}
\Fclass{\chi}{N,d} := \overline{\operatorname{span}}
\Big\{\textstyle\prod_{n\ge1}\big(\kappa^{(n)}_k\big)^{\otimes q_n}:
\textstyle\sum_n n\,q_n\le d\Big\},
\label{eq:Fclass}
\end{equation}

\begin{definition}[Connected order]\label{def:ordc}
For $F\in\Fclass{\chi}{N,d}$ written in the cumulant generators \eqref{eq:Fclass},
let $\ordc(F)$ be the largest cumulant order $n$ that appears as a factor in some
monomial of $F$. We call $F$ \emph{Gaussian-type} if $\ordc(F)\le 2$ and all its
order-1, 2 connected kernels are diagonal.
\end{definition}

\begin{lemma}[Gaussian closure]\label{lem:gauss}
Assume \ref{H1}--\ref{H2} and $\chi=0$. Then for every $n\ge3$ and every $p$,
\[
K^{(n)}_p\equiv0,
\]
i.e.\ the entire connected tower above order two vanishes identically; the state
is globally Gaussian. Moreover the surviving kernels are diagonal: for all
$p\ge1$,
\[
K^{(1)}_p(m_1,\dots,m_p)=\delta_{m_1\cdots m_p}\,S_1^{m_1-1}\beta^{(p)}_{(1)},
\qquad
K^{(2)}_p(m_1,\dots,m_p)=\delta_{m_1\cdots m_p}\,S_2^{m_1-1}\gamma^{(p)}_{(1)},
\]
where $\delta_{m_1\cdots m_p}=1$ iff $m_1=\cdots=m_p$ and
$\beta^{(p)}_{(1)},\gamma^{(p)}_{(1)}$ are the $p$-th Taylor coefficients of
$\beta,\gamma$ in the mode-1 block. Consequently every element of
$\Fclass{0}{N,d}$ is Gaussian-type in the sense of
Definition~\ref{def:ordc}: its only cross-time content is the disconnected (Wick)
product of diagonal order-1, 2 kernels.
\end{lemma}

\begin{lemma}[Kerr activation]\label{lem:kerr}
Assume \ref{H1}--\ref{H4} and $\chi\neq0$. Then:
\begin{enumerate}[label=(\alph*)]
\item the first moment acquires a connected cross-time bilinear kernel,
\[
K^{(1)}_2(m,m') = \chi\sum_{\ell\ge0}\big[e^{A_1\delta}\big]_\ell\,
\Phi_{m-\ell-\ell_\tau}\,\Psi_{m'-\ell-\ell_\tau} + (m\leftrightarrow m')\neq0
\quad (m\neq m'),
\]
where $\Phi_m:=S_1^{m-1}\beta_1$ is the mean response basis and
$\Psi_{m'}:=S_2^{m'-1}\gamma_1$ the covariance response basis;
\item the connected third cumulant is sourced,
\[
\kappa^{(3)}(t) = \chi\int_{-\infty}^t e^{A_3(t-s)}
V_3\!\big[\kappa^{(2)},\kappa^{(1)}\big](s-\tau)\,ds + O(\chi^2)\not\equiv0,
\]
so $K^{(3)}_p\not\equiv0$ for some $p$.
\end{enumerate}
\end{lemma}

\begin{corollary}[Unconditional $d=1$ witness]\label{cor:witness}
Let the target be the bilinear memory functional $y_k=u_{k-1}u_{k-3}$ and let the
readout degree be $d=1$ (linear in the homodyne quadratures). Then under
\ref{H1}--\ref{H4}:
\begin{enumerate}[label=(\roman*)]
\item for $\chi=0$ the achievable mean-square error is bounded below by a
strictly positive constant $\varepsilon_0(\mu)>0$ for every choice of $W$ and
every encoding; while
\item for $\chi\neq0$ the target lies in the closure $\Fclass{\Kerr}{N,1}$ and the
error can be driven below any $\varepsilon>0$.
\end{enumerate}
\end{corollary}

\begin{theorem}[Mode-count efficiency, general $\nu$]
\label{thm:scaling}
Assume \ref{H1}--\ref{H4}.
\begin{enumerate}[label=(\alph*)]
\item \textbf{(Gaussian confinement at fixed degree.)} For $\chi=0$ and readout
degree $d$, an order-$\nu$ connected kernel exists only for $d\ge\nu$ and is
assembled from $a$ mean factors and $b$ covariance factors with $a+b=\nu$ and
degree budget $a+2b\le d$, hence $b\le d-\nu$. Each temporal unfolding has rank at
most the dimension of the response space available to that slot: $2N$ for a
mean slot, $N(2N{+}1)$ for a covariance slot. In particular at the minimal degree
$d=\nu$ one has $b=0$: every slot is mean-confined and every unfolding has rank
$\le 2N$, independent of $M$.
\item \textbf{(Kerr enrichment.)} For $\chi\neq0$ the NM-Kerr reservoir realizes
order-$\nu$ connected kernels at readout degree $\le2$
(Proposition~\ref{prop:degree}) whose temporal unfoldings have rank growing with
the \emph{non-Markovian feedback depth} $D$---the number of independent delay
round-trips the feedback sustains above the loss floor---attaining rank
$\min(D,M)$ at fixed $N$. A single mode's feedback supplies $D$ distinct delay
channels, each carrying its own accumulated Kerr phase, so the connected kernel is
non-separable of rank up to $D$; this exceeds the Gaussian ceiling $2N$ whenever
$D>2N$.
\item \textbf{(Resource gap, two regimes.)} To realize a connected kernel of rank
$r$ (set by the non-Markovian feedback depth, $r\le\min(D,M)$) at the minimal
readout degree $d=\nu$, the Gaussian reservoir is mean-confined and requires
$2N\ge r$, i.e.\ $N_G=\Theta(r)$ modes; with unbounded readout degree it may use
the covariance family and requires $N(2N{+}1)\ge r$, i.e.\ $N_G=\Theta(\sqrt r)$.
The NM-Kerr reservoir attains it at $N_K=O(1)$ with feedback depth $D\ge r$. Since
reading high-degree quadrature polynomials is the dominant experimental cost, the
fixed-degree regime ($N_G=\Theta(r)$) is the operative one. The separating
resource is the feedback depth $D$: it is supplied by a single nonlinear mode and
matched by the Gaussian reservoir only with $\Theta(D)$ modes.
\end{enumerate}
\end{theorem}

\begin{proposition}[Degree folding]
\label{prop:degree}
Assume \ref{H1}--\ref{H4} and $\chi\neq0$. A connected order-$\nu$ functional is
accessible to the NM-Kerr reservoir at readout degree
\[
d_K = \begin{cases} 1, & \nu\ \text{odd},\\ 2, & \nu\ \text{even},\end{cases}
\qquad\text{at order }\chi^{\,q},\ q=\big\lceil(\nu-1)/2\big\rceil,
\]
whereas the Gaussian reservoir requires readout degree $d_G=\nu$.
\end{proposition}

\begin{proof}
\emph{Gaussian side.} By Lemma~\ref{lem:gauss} a connected order-$\nu$ functional
is reachable only as a disconnected product of $\nu$ single-time factors
(part (a) of the proof of Theorem~\ref{thm:scaling}); a $\nu$-fold product is a
degree-$\nu$ feature, so $d_G=\nu$ is necessary.

\emph{Kerr side.} The Kerr Hamiltonian $H_K=\tfrac{\chi}{2}a^{\dagger2}a^2$ gives,
through $\dot a=-i\chi a^\dagger a a$ and $a=(q+ip)/\sqrt2$, the symmetrized
quadrature drift
\begin{equation}
\dot q = \tfrac{\chi}{2}\big(p^3+pq^2\big),\qquad
\dot p = -\tfrac{\chi}{2}\big(q^3+qp^2\big),
\label{eq:kerr-drift}
\end{equation}
which is \emph{cubic}. Consequently the equation of motion for any order-$n$
symmetric moment, $\frac{d}{dt}\langle X_1\cdots X_n\rangle
=\sum_i\langle X_1\cdots\dot X_i\cdots X_n\rangle$, contains a contribution in
which the cubic drift raises the order by two, injecting an order-$(n{+}2)$
cumulant into the order-$n$ equation. Expanding in cumulants, the first-moment
equation reads
\[
\dot{\bar q}=\tfrac{\chi}{2}\big(K_{ppp}+K_{qqp}\big)+(\text{lower-order terms}),
\]
so a connected third-order cumulant sources $\kappa^{(1)}$ at order $\chi$ with
coefficient $\tfrac{\chi}{2}$, \emph{independent of the mean} $\bar q$. The
analogous computation for the third-moment equation gives
$\dot K_{qqq}\supset\tfrac{3\chi}{2}K_{qqppp}+\cdots$, folding order five into
order three with coefficient $\tfrac{3\chi}{2}$. By induction, the order-$n$
equation receives the order-$(n{+}2)$ cumulant with coefficient
$n\cdot\tfrac{\chi}{2}\cdot c$, $c\in\mathbb{Z}_{>0}$; the contribution never
cancels because the all-$p$ monomial $p^3$ in \eqref{eq:kerr-drift} always supplies
the all-$p$ cumulant term, which has no partner to cancel against. Iterating
$q=\lceil(\nu-1)/2\rceil$ times carries a connected order-$\nu$ functional into
$\kappa^{(1)}$ ($\nu$ odd) or $\kappa^{(2)}$ ($\nu$ even) with a nonzero
coefficient $\propto\chi^{q}$. A functional in $\kappa^{(1)}$ (resp.\
$\kappa^{(2)}$) is a degree-1 (resp.\ degree-2) homodyne feature, so $d_K\in\{1,2\}$.
The base case $\nu=2$ is Lemma~\ref{lem:kerr}(a)/Corollary~\ref{cor:witness}.
\end{proof}

\end{document}
